\begin{frame}{SAR: Aproximações Sucessivas}
\begin{enumerate}
\item O sample-and-hold mantém a tensão de entrada durante a conversão.
\item O SAR testa um bit por vez, do mais significativo ao menos significativo. O DAC interno gera a tensão de teste.
\item O comparador decide: se a tensão de teste ultrapassar a entrada, o bit volta a zero. Caso contrário, permanece em um.
\end{enumerate}
\begin{block}{Exemplo didático com 4 bits}
Para $V_{in}=1{,}50$ V e $V_{ref}=2{,}56$ V:
\begin{center}
\small
\begin{tabular}{ccc}
\textbf{Código testado} & \textbf{Tensão do DAC} & \textbf{Decisão}\\\hline
\texttt{1000} & 1,28 V & Mantém o bit\\
\texttt{1100} & 1,92 V & Retira o bit\\
\texttt{1010} & 1,60 V & Retira o bit\\
\texttt{1001} & 1,44 V & Mantém o bit\\
\end{tabular}
\end{center}
\end{block}
\pblfonte{Datasheet, seção 26.2. Exemplo simplificado com comparador e DAC ideais.}
\end{frame}

\begin{frame}{Por que Utilizamos o SAR?}
\begin{itemize}
\item \textbf{Já faz parte do ATmega2560:} utilizamos o ADC integrado à plataforma definida no problema.
\item \textbf{Atende à velocidade dos ensaios:} a conversão normal leva 13 ciclos do ADC, ou $13\ \mu$s com clock de 1 MHz.
\item \textbf{Permite controlar a aquisição:} configuramos referência, canal e prescaler pelos registradores. O Timer1 define o intervalo entre amostras.
\item \textbf{Entrega a resolução usada no projeto:} lemos os 8 bits mais significativos para representar 256 níveis e transmitir uma amostra por byte.
\end{itemize}
\begin{alertblock}{Leitura de 8 bits}
O ADC continua realizando a conversão interna de 10 bits. Ler apenas \texttt{ADCH} não transforma o hardware em um SAR de 8 etapas.
\end{alertblock}
\pblfonte{Datasheet, seções 26.2 e 26.4; configuração em adc\_uart.ino.}
\end{frame}
